\section{Introduction}

This thesis was written on the topic \emph{Development of a domain expert using a large language model like ChatGPT}. Large language models (LLMs) of the GPT-4 class answer many medical questions at the level of a well-read student \cite{singhal2023}, so it is tempting to use such a model directly as an ``expert''. In a clinical field, however, a general-purpose model has three weaknesses: it can produce fluent but unsupported claims (hallucination \cite{ji2023}); it offers no provenance, and may even invent references that do not exist \cite{walters2023}; and it has no access to the case --- here, the cells on a blood-smear image --- so its answers are generic. In this thesis, a \emph{domain expert} is therefore an LLM-based system that answers within one field, grounds its statements in curated domain literature and cites it, inspects case-specific evidence through tools, and signals when a human expert must take over.

\paragraph{Choice of application.} The topic is addressed through a concrete system: a \emph{hybrid multimodal lab assistant} for peripheral blood smears. The blood smear is one of the most informative tests in hematology: a stained drop of blood is examined under the microscope, and the proportions of the white-cell types (the \emph{differential}) can reveal infection, leukaemia or bone-marrow disease. Manual review is slow and observer-dependent, which has motivated much work on automated cell classification \cite{acevedo2019,habibzadeh2018}; but such systems stop at a label per cell. Turning the labels into a clinical interpretation requires domain knowledge. Hematology is therefore a natural test bed for an LLM domain expert, because it combines visual, case-specific evidence, which image models can extract, with textbook knowledge, which the expert must supply.

\paragraph{Theoretical and practical importance.} Theoretically, the thesis compares the main ways of turning a general LLM into a domain expert --- retrieval-augmented generation (RAG), tool use and parameter-efficient fine-tuning --- on the same domain, and asks what each of them actually contributes. Practically, it shows how such an expert can be connected to image analysis in a way that is inspectable and safe: every automated finding carries an uncertainty estimate, every claim can be traced to a retrieved source, and the decision to involve a human is taken by explicit rules rather than by the model. It also examines whether an open model that can run on local hardware can replace a commercial API, which matters wherever patient data may not leave the institution.

\subsection{Aims and research questions}

The aim is to build and evaluate the lab assistant and, through it, to answer four research questions:

\begin{description}[leftmargin=1.2cm,style=nextline]
  \item[RQ1] Can standard deep-learning models (YOLOv8, EfficientNet-B0) detect and classify blood cells accurately and attach a useful uncertainty estimate and a saliency map to every classified white cell?
  \item[RQ2] How can a general-purpose LLM be turned into a hematology domain expert with retrieval-augmented generation and domain-specific tools, and what does the resulting assistant produce on real smear images?
  \item[RQ3] Does QLoRA fine-tuning of an open-weight model (Llama-3.1-8B-Instruct) on textbook-grounded question--answer pairs improve its hematology answers compared with the same model without fine-tuning?
  \item[RQ4] Can the combined system make its reasoning inspectable and hand difficult cases to a human through deterministic rules --- and where does this fail?
\end{description}

\subsection{Research methods}

The work combines system development with quantitative evaluation. The image-analysis stages were trained on public benchmark datasets (TXL-PBC for detection, PBC for classification) and evaluated on held-out test sets with standard metrics, complemented by a calibration and error-detection analysis of the uncertainty estimate. The domain expert was built from a knowledge base of 1{,}049 textbook passages, a GPT-4o agent with five hematology tools, and a Llama-3.1-8B model fine-tuned with QLoRA on synthetic question--answer pairs generated from the same textbooks. The fine-tuned model was compared with its base version and with GPT-4o on held-out questions (similarity to reference answers and ratings by an LLM judge) and on external exam questions, using paired statistical tests. Finally, a complete run of the assistant was analysed as a case study. All code, notebooks, configuration and results are part of the submitted software artefact.

The main findings, detailed in Chapter~\ref{sec:results}, are that the image models perform strongly on the benchmark data (98.60\,\% classification accuracy, with uncertainty scores that detect errors with an AUROC of 0.97), and that fine-tuning changed \emph{how} the open model answers --- closer to the textbook, without invented citations --- but not \emph{what} it knows: its clinical correctness did not improve, and GPT-4o remained the most correct system. Knowledge and provenance therefore have to come from retrieval, and safety from explicit rules.

\subsection{Structure of the thesis}

Chapter~\ref{sec:background} reviews how domain-expert LLMs are built and the related work on blood-smear analysis, uncertainty and interpretability. Chapter~\ref{sec:data} describes the data and knowledge sources, and Chapter~\ref{sec:arch} the architecture of the lab assistant. Chapter~\ref{sec:method} gives the training procedures and the evaluation protocol, and Chapter~\ref{sec:results} presents and discusses the results, including the case study and the limitations. Chapter~\ref{sec:conclusion} concludes and outlines future work.

\clearpage
